% ============================================================
%  preamble.tex — packages and settings shared by the whole thesis
%  Compile with XELATEX (not pdflatex). On Overleaf: Menu -> Compiler -> XeLaTeX
%  Why XeLaTeX: the thesis text contains many raw Unicode symbols
%  (σ, µ, °, ±, ≥, ≤, →, ↑, ↓, √, Δ, ₀) copied from the drafted chapters.
%  XeLaTeX renders any Unicode codepoint the loaded font supports, with no
%  inputenc/fontenc fiddling. pdfLaTeX will choke on several of these.
% ============================================================

% ---- Page geometry -------------------------------------------------------
\usepackage[a4paper,margin=1in,bindingoffset=0.5in]{geometry}

% ---- Fonts (XeLaTeX) ------------------------------------------------------
\usepackage{fontspec}
\setmainfont{TeX Gyre Termes}      % wide Unicode coverage, Times-like
\setsansfont{TeX Gyre Heros}
\setmonofont{DejaVu Sans Mono}[Scale=0.85]
% DejaVu Sans Mono (not TeX Gyre Cursor) specifically because the code
% listings in Appendix E contain box-drawing divider comments
% (# ── Hyper-parameters ── ...) copied verbatim from the real source files;
% TeX Gyre Cursor is missing those glyphs, DejaVu covers them.
% If TeX Gyre fonts are unavailable in your Overleaf/local install, also swap:
%   \setmainfont{DejaVu Serif}    (very broad symbol coverage, safest fallback)

% ---- Maths -----------------------------------------------------------------
\usepackage{amsmath,amssymb}

% ---- Arabic front-matter support ------------------------------------------
% DEFAULT APPROACH: include a pre-rendered PDF page for the Arabic abstract
% (see frontmatter/abstract_ar.tex and frontmatter/abstract_ar.pdf). This is
% robust and needs nothing else in this preamble.
\usepackage{pdfpages}
%
% ALTERNATIVE: native Arabic via polyglossia+bidi was tried for this skeleton
% and abandoned deliberately. The `bidi` package (loaded by polyglossia the
% moment \setotherlanguage{arabic} is used) imposes a strict, semi-documented
% load order on OTHER packages in this preamble (amsmath and geometry must
% load BEFORE bidi; caption, hyperref, csquotes must load AFTER it — and the
% exact set varies by bidi version). Getting every package below into the
% order bidi wants is fiddly and version-fragile, which is a bad trade for a
% single abstract page on a first LaTeX project. If you want native Arabic
% text later (e.g. to also translate other sections), uncomment the four
% lines below, then re-order the *rest of this file* so every package that
% conflicts loads on the correct side of \usepackage{polyglossia} — compile,
% read the "load X before/after bidi" error, move that one package, repeat.
%
% \usepackage{polyglossia}
% \setmainlanguage{english}
% \setotherlanguage{arabic}
% \newfontfamily\arabicfont[Script=Arabic]{Amiri}

% ---- Tables ------------------------------------------------------------
\usepackage{booktabs}      % \toprule \midrule \bottomrule — used by pandoc's table output
\usepackage{longtable}     % tables that break across pages (Appendix A, D)
\usepackage{array}
\usepackage{multirow}
\usepackage{tabularx}

% ---- Figures ---------------------------------------------------------------
\usepackage{graphicx}
\usepackage{subcaption}
\usepackage{float}

% ---- Code listings (Appendix E) --------------------------------------------
\usepackage{listings}
\usepackage{xcolor}
\lstdefinestyle{pythonstyle}{
  language=Python,
  basicstyle=\ttfamily\small,
  keywordstyle=\color{blue!70!black}\bfseries,
  commentstyle=\color{gray}\itshape,
  stringstyle=\color{orange!70!black},
  numbers=left, numberstyle=\tiny\color{gray}, numbersep=8pt,
  frame=single, breaklines=true, showstringspaces=false,
  tabsize=4, columns=flexible
}
\lstset{style=pythonstyle}
% (If you prefer real syntax-highlighting fidelity, switch to the `minted`
%  package instead — needs -shell-escape, which Overleaf enables by default
%  under Menu -> Settings, but `listings` above needs no extra setup.)

% ---- Cross-referencing ------------------------------------------------------
\usepackage{hyperref}
\hypersetup{
  colorlinks=true, linkcolor=blue!50!black, citecolor=blue!50!black,
  urlcolor=blue!50!black, bookmarksnumbered=true
}
\usepackage{cleveref}      % \cref{}, \Cref{} — use instead of \ref{} everywhere

% ---- Misc ------------------------------------------------------------------
\usepackage{csquotes}
\usepackage[nottoc]{tocbibind}   % puts Bibliography in the ToC automatically
\usepackage{fancyhdr}
\pagestyle{fancy}
\fancyhf{}
\fancyhead[LE,RO]{\thepage}
\fancyhead[RE]{\nouppercase{\leftmark}}
\fancyhead[LO]{\nouppercase{\rightmark}}
\renewcommand{\headrulewidth}{0.4pt}
\setlength{\headheight}{15pt}   % silences fancyhdr's "headheight too small" warning

% ---- Pandoc compatibility --------------------------------------------------
% Pandoc emits \tightlist inside compact (no-blank-line) markdown lists.
% It is not a standard LaTeX command, so every converted chapter needs this.
\providecommand{\tightlist}{%
  \setlength{\itemsep}{0pt}\setlength{\parskip}{0pt}}

% ---- Custom commands used across chapters -----------------------------------
\newcommand{\rqref}[1]{\textbf{RQ#1}}       % \rqref{4} -> RQ4  (named rqref: \rq is a LaTeX primitive)
\newcommand{\code}[1]{\texttt{#1}}          % inline code / filenames
\newcommand{\pct}[1]{#1\%}                  % percentages, keeps % escaping consistent

% ---- Table caption style ----------------------------------------------------
\usepackage{caption}
\captionsetup{font=small,labelfont=bf,skip=6pt}
